
\documentclass{article}
\usepackage{anysize} 
\usepackage{multirow}
\usepackage{float}
\usepackage{booktabs}
\usepackage{graphicx}
\begin{document}



\section{PART I}

\subsection{Cautionary Statement About Forward-Looking Statements}
This Annual Report on Form 10-K contains certain forward-looking statements that are made subject to the safe harbor protections provided by the Private Securities Litigation Reform Act of 1995. Forward-looking statements are often identified by the words, “will,” “may,” “should,” “continue,” “anticipate,” “believe,” “expect,” “target,” “plan,” “forecast,” “project,” “estimate,” “intend,” “commit,” “potential” and words of a similar nature and generally include statements regarding future results of operations, including revenues, earnings or cash flows; plans and objectives for the future; projections, estimates or assumptions relating to our operational or financial performance, including anticipated impacts of the Inflation Reduction Act of 2022; projections, estimates or assumptions relating to our capital expenditures; or opinions, views or beliefs about the effects of current or future events, circumstances or performance.

You should view these statements with caution. These statements are not guarantees of future performance, circumstances or events. They are based on facts and circumstances known to us as of the date the statements are made, and you should not place undue reliance on any such forward-looking statements. Forward-looking statements are subject to certain risks and uncertainties that could cause actual results to differ materially from our historical experience and our present expectations or anticipated results. These risks and uncertainties include, but are not limited to, those described in Part 1, “Item 1A. Risk Factors” and elsewhere in this report and may also be described from time to time in our future reports filed with the U.S. Securities and Exchange Commission (“SEC”). We do not undertake any obligation to update forward-looking statements to reflect events, circumstances, changes in expectations or other developments after the date of those statements.

\subsection{Item 1. Business}

\subsubsection{General}
Waste Management, Inc. is a holding company and all operations are conducted by its subsidiaries. When the terms “the Company,” “we,” “us” or “our” are used in this document, those terms refer to Waste Management, Inc., together with its consolidated subsidiaries and consolidated variable interest entities. When we use the term “WMI,” we are referring only to Waste Management, Inc., the parent holding company.

WMI was incorporated in Oklahoma in 1987 under the name “USA Waste Services, Inc.” and was reincorporated as a Delaware company in 1995. In a 1998 merger, the Illinois-based waste services company formerly known as Waste Management, Inc. became a wholly-owned subsidiary of WMI and changed its name to Waste Management Holdings, Inc. (“WM Holdings”). At the same time, our parent holding company changed its name from USA Waste Services to Waste Management, Inc. Like WMI, WM Holdings is a holding company and all operations are conducted by subsidiaries.

Our principal executive offices are located at 800 Capitol Street, Suite 3000, Houston, Texas 77002. Our telephone number is (713) 512-6200. Our website address is www.wm.com. Our annual reports on Form 10-K, quarterly reports on Form 10-Q and current reports on Form 8-K are all available, free of charge, on our website as soon as practicable after we file the reports with the SEC. Our stock is traded on the New York Stock Exchange under the symbol “WM.”

We are North America’s leading provider of comprehensive environmental solutions, providing services throughout the United States (“U.S.”) and Canada. We partner with our customers and the communities we serve to manage and reduce waste at each stage from collection to disposal, while recovering valuable resources and creating clean, renewable energy. Our solid waste business is operated and managed locally by our subsidiaries that focus on distinct geographic areas and provide collection, transfer, disposal, recycling and resource recovery services. Through our subsidiaries, including our Waste Management Renewable Energy (“WM Renewable Energy”) business, we are also a leading developer, operator and owner of landfill gas-to-energy facilities in the U.S. and Canada that produce renewable electricity and renewable natural gas, which is a significant source of fuel that we allocate to our natural gas fleet. During 2023, our largest customer represented less than 5\% of annual revenues.

We own or operate 263 landfill sites, which is the largest network of landfills throughout the U.S. and Canada. In order to make disposal more practical for larger urban markets, where the distance to landfills is typically farther, we manage 332 transfer stations that consolidate, compact and transport waste efficiently and economically. We also use waste to create energy, recovering the gas produced naturally as waste decomposes in landfills and using the gas in generators to make electricity. We are a leading recycler in the U.S. and Canada, handling materials that include cardboard, paper, glass, plastic and metal. We provide cost-efficient, environmentally sound recycling programs for municipalities, businesses and households across the U.S. and Canada as well as other services that supplement our solid waste business.

Our fundamental strategy has not changed; we remain dedicated to providing long-term value to our stockholders by successfully executing our core strategy of focused differentiation and continuous improvement. We have enabled a people-first, technology-led focus to drive our mission to maximize resource value, while minimizing environmental impact, and sustainability and environmental stewardship is embedded in all that we do. Our strategy leverages and sustains the strongest asset network in the industry to drive best-in-class customer experience and growth. Our strategic planning processes appropriately consider that the future of our business and the industry can be influenced by changes in economic conditions, the competitive landscape, the regulatory environment, asset and resource availability and technology. We believe that focused differentiation, which is driven by capitalizing on our unique and extensive network of assets, will deliver profitable growth and position us to leverage competitive advantages. Simultaneously, we believe that investing in automation to improve processes and drive operational efficiency combined with a focus on the cost to serve our customer will yield an attractive profit margin and enhanced service quality. We are furthering our strategy of focused differentiation and continuous improvement beyond our traditional waste operations through our sustainability growth strategy that includes significant planned investments in our WM Renewable Energy and Recycling Processing and Sales businesses, while increasing automation and reducing labor dependency. We are also evaluating and pursuing emerging diversion technologies that may generate additional value.

Our Company's goals are targeted at putting our people first, positioning them to serve and care for our customers, the environment, the communities in which we work and our stockholders. Our brand promise is ALWAYS WORKING FOR A SUSTAINABLE TOMORROW®. We live this promise through our service offerings and sustainable solutions, our investments in innovation, our people, and our commitment to the future. Through our longtime focus on finding sustainable solutions, we continue to evolve beyond being a traditional environmental waste services company. Increasingly, our industry-leading focus on environmental sustainability aligns with demand from our customers who want more of their waste materials recovered. Waste streams are becoming more complex, and our aim is to address current needs, while anticipating the expanding and evolving needs of our customers. We believe we are uniquely equipped to meet the challenges of the changing waste industry and our customers’ waste management needs, both today and tomorrow as we work together to envision and create a more sustainable future.

We believe that execution of our strategy will deliver shareholder value and leadership in a dynamic industry and in any economic environment. In addition, we intend to continue to return value to our stockholders through dividend payments and our common stock repurchase program. In December 2023, we announced that our Board of Directors expects to increase the quarterly dividend from \$0.70 to \$0.75 per share for dividends declared in 2024, which is a 7.1\% increase from the quarterly dividends we declared in 2023. This is an indication of our ability to generate strong and consistent cash flows and marks the 21st consecutive year of dividend increases. All quarterly dividends will be declared at the discretion of our Board of Directors and depend on various factors, including our net earnings, financial condition, cash required for future business plans, growth and acquisitions and other factors the Board of Directors may deem relevant.

\subsubsection{Operations}

\paragraph{General}

To enhance transparency regarding our financial performance, highlight the strength and consistency of our core solid waste businesses, and underscore our commitment to sustainability through planned and ongoing investments in our Recycling Processing and Sales, and WM Renewable Energy businesses, beginning in the fourth quarter of 2023, our senior management revised its segment reporting to (i) reflect the financial results of our collection, transfer, disposal and resource recovery services businesses independently; (ii) combine the results of all recycling facilities from our East and West Tier segments with our recycling brokerage and sales activities to form a newly created Recycling Processing and

Sales reportable segment and (iii) include our WM Renewable Energy business as a reportable segment. Accordingly, our senior management now evaluates, oversees and manages the financial performance of our business through four reportable segments, referred to as (i) Collection and Disposal - East Tier ("East Tier"); (ii) Collection and Disposal - West Tier ("West Tier"); (iii) Recycling Processing and Sales; and (iv) WM Renewable Energy. Our East and West Tiers, along with certain ancillary services ("Other Ancillary") not managed through our Tier segments, but that support our collection and disposal operations, form our "Collection and Disposal" businesses.

Our East Tier primarily consists of geographic areas located in the Eastern U.S., the Great Lakes region and substantially all of Canada. Our West Tier primarily includes geographic areas located in the Western U.S., including the upper Midwest region, and British Columbia, Canada.

We also provide additional services not managed through our four reportable segments, which are presented as Corporate and Other. For further discussion refer to Note 19 of our Consolidated Financial Statements. Reclassifications have been made to our prior period consolidated financial information to conform to the current year presentation.

\paragraph{Collection and Disposal}

Services provided through our Collection and Disposal businesses are described below:

\subparagraph{Collection}
Our commitment to customers begins with a vast waste collection network. Collection involves picking up and transporting waste and recyclable materials from where it was generated to a transfer station, recycling facility or disposal site. We generally provide collection services under one of two types of arrangements:

\begin{itemize}
    \item For commercial and industrial collection services, typically we have three-year service agreements. The fees under the agreements are influenced by factors such as collection frequency, type of collection equipment we furnish, type and volume or weight of the waste collected, distance to the disposal facility, labor costs, cost of disposal and general market factors. As part of the service, we provide steel containers to most customers to store their solid waste between pickup dates. Containers vary in size and type according to the needs of our customers and the restrictions of their communities. Many are designed to be lifted mechanically and are emptied into a truck’s compaction hopper or directly into a disposal site. By using these containers, we can service most of our commercial and industrial customers with trucks operated by only one employee.
    
    \item For most residential collection services, we have a contract with, or a franchise granted by, a municipality, homeowners’ association or some other regional authority that gives us the exclusive right to service all or a portion of the homes in an area. These contracts or franchises are typically for periods of three to ten years and typically mirror maximum terms as allowed by statutes by state. We also provide services directly to individual monthly subscriptions to households. The fees for residential collection are either paid by the municipality or authority from their tax revenues or service charges, or are paid directly by the residents receiving the service. The Company is generally phasing out traditional manual systems and moving to further automate residential collection services. Benefits of automation include enhanced worker safety, improved service delivery to the customer and an overall reduction in the cost to provide services.
\end{itemize}

\subparagraph{Landfill}
Landfills are the main depositories for solid waste in North America. As of December 31, 2023, we owned or operated 258 solid waste landfills and five secure hazardous waste landfills, which represents the largest network of landfills throughout the U.S. and Canada. As of December 31, 2023, we owned or controlled the management of 237 sites with remedial activities, that are in closure or that have received a certification of closure from the applicable regulatory agency. Solid waste landfills are constructed and operated on land with engineering safeguards that limit the possibility of water and air pollution, and are operated under procedures prescribed by regulation. A landfill must meet federal, state or provincial, and local regulations during its design, construction, operation and closure. The operation and closure activities of a solid waste landfill include excavation, construction of liners, continuous spreading and compacting of waste, covering of waste with earth or other acceptable material and constructing final capping of the landfill. These operations are carefully planned to maintain environmentally safe conditions and to maximize the use of the airspace.

All solid waste management companies must have access to a disposal facility, such as a solid waste landfill. The significant capital requirements of developing and operating a landfill serve as a barrier to landfill ownership and, thus,


\end{document}